\documentclass[11pt]{article}
\usepackage[margin=1in]{geometry}
\usepackage{amsmath,amssymb,amsthm,booktabs,graphicx,microtype,natbib}
\usepackage[hidelinks]{hyperref}
\usepackage{xcolor}
\usepackage{enumitem}

\newtheorem{theorem}{Theorem}
\newtheorem{proposition}{Proposition}
\newtheorem{corollary}{Corollary}
\newtheorem{remark}{Remark}
\newcommand{\E}{\mathbb E}
\newcommand{\Pp}{\mathbb P}
\newcommand{\Var}{\operatorname{Var}}
\newcommand{\1}{\mathbb I}
\newcommand{\F}{\mathcal F}
\newcommand{\given}{\,\middle|\,}

\title{When Nothing Happens:\\Martingale-Coherent Distribution Forecasts for Prediction Markets}
\author{Anonymous research artifact}
\date{July 2026}

\begin{document}
\maketitle

\begin{abstract}
Prediction-market prices often remain unchanged for many observation periods and
then move abruptly. Existing structural volatility models explain the size of
active moves but do not provide an unconditional predictive distribution for
the joint incidence and magnitude of price updates. We derive a restriction
that couples these two margins. If a binary-claim price is a bounded martingale,
the fraction of remaining Bernoulli uncertainty released over one period,
$r_t=\Var(p_{t+1}\mid\F_t)/[p_t(1-p_t)]$, cannot exceed the update probability
$q_t=\Pp(p_{t+1}\ne p_t\mid\F_t)$. We construct a martingale-coherent
hurdle-beta distribution that realizes every feasible pair $0<r_t<q_t\le1$,
contains a point mass at no change, respects the $[0,1]$ price boundary, and has
closed-form quantiles. We also replace the locally additive DR-AS variance with
a finite-step information-clock aggregation that cannot spend more than the
remaining uncertainty and agrees with DR-AS to first order. A public-API
replication evaluates the model in monthly expanding windows using proper
mixed-distribution scores, multi-level coverage, and contract-cluster
inference. Empirical results are reported only from the locked pipeline and are
not asserted in this draft until the full backtest and robustness suite finish.
\end{abstract}

\section{Introduction}

A prediction-market price is both a traded price and a probability. This dual
role makes the standard volatility question unusually structured. The price is
bounded, its terminal payoff is binary, and its resolution date is known. At
the same time, the observed quote can remain exactly unchanged for hours even
while uncertainty persists. A useful forecast must therefore answer two
questions: will the displayed probability update, and how far can it move if it
does?

\citet{xi2026volatility} derive a structural variance model with a
Wright-Fisher deadline-resolution channel and a Glosten-Milgrom order-flow
channel. Their DR-AS specification outperforms plain ARCH and GARCH models on a
large Kalshi panel, especially after residual GARCH dynamics are added. The
headline exercise deliberately conditions on active price updates. The paper's
conclusion identifies a joint update-hazard and move-size model as a natural
next step.

This paper develops that step as a distributional rather than merely
regression extension. The central observation is that update incidence and
conditional variance cannot be chosen independently once the price is required
to remain a bounded martingale. Let $q_t$ be the probability of any change and
let $r_t$ be unconditional one-period variance divided by remaining Bernoulli
uncertainty. We prove $r_t\le q_t$. The inequality is sharp, easy to test, and
economically interpretable: rare updates can release only a correspondingly
limited share of uncertainty unless an active update is nearly a full
resolution to zero or one.

The paper makes four contributions.

\begin{enumerate}[leftmargin=*]
\item It gives a necessary and sharp coherence restriction linking the
extensive margin of information arrival to the intensive margin of belief
revision.
\item It constructs a hurdle-beta predictive law with exact martingale mean,
exact unconditional variance, an explicit no-update atom, bounded support, and
closed-form central intervals.
\item It derives a finite-step information-clock version of DR-AS. The pure
deadline component in fact has the exact normalized variance $\Delta/\tau$;
order-flow intensity enters through an exponential clock map that preserves the
total variance budget and nests additive DR-AS locally.
\item It evaluates full predictive distributions on the unfiltered hourly
panel using scores that are proper for mixed and price-lattice observations,
rather than evaluating a variance proxy only after zero moves are removed.
\end{enumerate}

The construction differs from general zero-inflated price models
\citep{catania2019dynamic,holy2023intraday} because its two margins are coupled
by the martingale and binary-payoff structure. It differs from logit
jump-diffusion kernels for prediction markets \citep{dalen2025toward} because it
targets the observed no-update atom and an unconditional short-horizon density,
not a continuous-time derivative-pricing surface. Recent local-volatility work
for event-market martingales \citep{xu2026local} characterizes admissible
boundary behavior and information clocks; the result here instead identifies a
discrete-observation inequality between update incidence and variance release.

\section{Structural benchmark and forecasting target}

Let $\Theta\in\{0,1\}$ be the settlement payoff at time $T$. Under the maintained
risk-neutral interpretation,
\begin{equation}
p_t=\E[\Theta\mid\F_t]
\end{equation}
is a bounded martingale with $p_T=\Theta$. For an observation step of length
$\Delta$ and time to resolution $\tau=T-t$, the DR-AS benchmark of
\citet{xi2026volatility} is
\begin{equation}
h^2_{\mathrm{DR-AS}}
=\left[\frac{p_t(1-p_t)}{\tau}
+K\nu(V_t)\frac{s_t^2}{4}\right]\Delta,
\label{eq:dras}
\end{equation}
where $s_t$ is the bid-ask spread, $V_t$ is volume, $\nu$ is a concave activity
transform, and $K\ge0$ is fitted.

The benchmark maps $h$ to a symmetric normal-reference interval around $p_t$
and clips that interval to $[0,1]$. This provides a common scoring rule for scale
forecasts but is not the transition law of either structural channel. More
importantly for the present paper, conditioning estimation and headline
evaluation on $p_{t+\Delta}\ne p_t$ changes the variance object. We instead
forecast the full law of $p_{t+\Delta}$ on the unfiltered panel.

\section{A coherence law for update incidence}

Fix a forecast origin and suppress time subscripts. Write $p\in(0,1)$ for the
current price, $Y\in[0,1]$ for the next price, and
\begin{equation}
q=\Pp(Y\ne p\mid\F),\qquad
r=\frac{\Var(Y\mid\F)}{p(1-p)}.
\end{equation}

\begin{theorem}[Update-variance coherence]
\label{thm:coherence}
If $Y\in[0,1]$ almost surely and $\E[Y\mid\F]=p$, then
\begin{equation}
0\le r\le q\le1.
\label{eq:coherence}
\end{equation}
If $q>0$, equality $r=q$ holds if and only if the active-update distribution is
supported on $\{0,1\}$ almost surely. If $r=0$, then $q=0$ under the definition
$q=\Pp(Y\ne p\mid\F)$.
\end{theorem}

\begin{proof}
Let $A=\{Y\ne p\}$. Since $Y-p=0$ on $A^c$ and $\E[Y-p\mid\F]=0$,
$\E[Y\mid A,\F]=p$ whenever $q>0$. Hence
\begin{align}
\Var(Y\mid\F)
&=q\E[(Y-p)^2\mid A,\F]\\
&=q\Var(Y\mid A,\F)\\
&\le q\,p(1-p),
\end{align}
where the last line is the Bernoulli variance bound for a random variable in
$[0,1]$ with mean $p$: $\E[Y^2\mid A,\F]\le\E[Y\mid A,\F]=p$. Equality in that
bound requires $Y^2=Y$ almost surely conditional on $A$, which is equivalent to
$Y\in\{0,1\}$. If $r=0$, then $Y=p$ almost surely, so $q=0$.
\end{proof}

The result does not rely on a diffusion, Gaussian shocks, or a particular
market microstructure model. It is a consequence of three economically basic
restrictions: bounded prices, the martingale mean, and an observed atom at no
change. It also supplies a diagnostic. A separately estimated update hazard
below a structural variance release cannot correspond to any bounded
martingale transition with those moments.

\section{Finite-step structural variance release}

The neutral Wright-Fisher diffusion in information time $u$ is
\begin{equation}
d\widetilde p_u=\sqrt{\widetilde p_u(1-\widetilde p_u)}\,dB_u.
\end{equation}
Its generator applied to $g(p)=p(1-p)$ gives $\mathcal Lg=-g$. Therefore
\begin{equation}
\E[\widetilde p_u(1-\widetilde p_u)\mid\widetilde p_0=p]
=p(1-p)e^{-u}
\end{equation}
and, using the martingale mean,
\begin{equation}
\Var(\widetilde p_u\mid\widetilde p_0=p)
=p(1-p)(1-e^{-u}).
\label{eq:wfexact}
\end{equation}

The deadline clock has instantaneous rate $1/(T-t)$. Over a calendar step
$[t,t+\Delta]$ with $\tau>\Delta$, its integrated information time is
\begin{equation}
\Lambda_{\mathrm{DR}}
=\int_t^{t+\Delta}\frac{ds}{T-s}
=\log\left(\frac{\tau}{\tau-\Delta}\right).
\end{equation}
Combining this expression with Equation~\eqref{eq:wfexact} yields
\begin{equation}
r_{\mathrm{DR}}=1-e^{-\Lambda_{\mathrm{DR}}}=\frac{\Delta}{\tau}.
\label{eq:drexact}
\end{equation}
Thus, under the maintained Wright-Fisher clock, the pure deadline variance in
Equation~\eqref{eq:dras} is exact for a finite step ending before the deadline,
even though it can also be obtained by an Euler argument.

To combine order flow without violating the remaining-uncertainty budget,
define its variance-equivalent local clock intensity
\begin{equation}
\lambda_{\mathrm{AS}}
=K\nu(V_t)\frac{s_t^2}{4p_t(1-p_t)}.
\end{equation}
We propose the finite-step release
\begin{equation}
r_t(K)
=1-\exp\left\{-\Lambda_{\mathrm{DR}}
-\lambda_{\mathrm{AS}}\Delta\right\}.
\label{eq:release}
\end{equation}
It lies in $[0,1]$ by construction. A small-step expansion gives
\begin{equation}
r_t(K)p_t(1-p_t)
=\left[\frac{p_t(1-p_t)}{\tau}
+K\nu(V_t)\frac{s_t^2}{4}\right]\Delta+O(\Delta^2),
\end{equation}
so the model nests DR-AS locally while introducing the finite-horizon
interaction required by a single variance budget. The order-flow clock is a
variance-equivalent reduced form; it is not claimed to turn the exact
Glosten-Milgrom jump distribution into a Wright-Fisher diffusion.

\section{A martingale-coherent hurdle distribution}

For any $0<r<q\le1$, set
\begin{equation}
\kappa=\frac{q}{r}-1
\end{equation}
and define
\begin{equation}
Y\mid\F\sim(1-q)\delta_p
+q\,\mathrm{Beta}\bigl(p\kappa,(1-p)\kappa\bigr).
\label{eq:mhb}
\end{equation}

\begin{proposition}[Sharp construction]
The law in Equation~\eqref{eq:mhb} has
\begin{equation}
\E[Y\mid\F]=p,\qquad
\Var(Y\mid\F)=r p(1-p),\qquad
\Pp(Y\ne p\mid\F)=q.
\end{equation}
For $q=r$, its weak endpoint limit replaces the active beta component by a
Bernoulli$(p)$ law on $\{0,1\}$. Together with the degenerate $r=q=0$ case, the
construction realizes every feasible pair in Theorem~\ref{thm:coherence}.
\end{proposition}

\begin{proof}
The active beta component has mean $p$ and variance
$p(1-p)/(\kappa+1)=(r/q)p(1-p)$. Mixture with a point mass at the same mean
preserves the mean and multiplies active variance by $q$. A continuous beta draw
equals $p$ with probability zero, giving update probability $q$. As
$\kappa\downarrow0$, the beta law converges weakly to mass $1-p$ at zero and
mass $p$ at one.
\end{proof}

The beta component is a parsimonious moment-matching kernel. It is not the exact
finite-time Wright-Fisher transition density, which has nontrivial absorption
behavior and is difficult to express analytically \citep{roa2022smalltime}.
This distinction is important: our theoretical claim concerns exact moments and
the update atom, not exact recovery of the diffusion transition law.

We parameterize the hazard so coherence holds at every forecast origin:
\begin{equation}
q_t=r_t+(1-r_t)\operatorname{logit}^{-1}(x_t'\gamma),
\label{eq:hazard}
\end{equation}
where $x_t$ contains only forecast-origin states: log volume, log open interest,
log time to resolution, log spread, $p_t(1-p_t)$, and the lagged update
indicator. Equation~\eqref{eq:hazard} can be read as a structural floor $r_t$
plus a residual probability of a partial, nonterminal quote update.

\section{Estimation and proper evaluation}

Each monthly test window is forecast using parameters estimated only on prior
calendar months. We fit $K$ on the full training panel by volume-weighted Gaussian
quasi-likelihood for the unconditional variance $r_t(K)p_t(1-p_t)$. Conditional
on this release, $\gamma$ is estimated by Bernoulli likelihood for the update
indicator under Equation~\eqref{eq:hazard}, using the same forecast-origin volume
weights. This two-stage procedure targets
the identifying moments directly and avoids using the future update indicator
as a forecast-origin regressor.

The primary comparison contains five models:
\begin{enumerate}[leftmargin=*]
\item deadline resolution with clipped normal-reference intervals;
\item the original active-update DR-AS fit with clipped normal-reference
intervals;
\item the bounded beta kernel with $q=1$, isolating distribution shape without
a hurdle;
\item a separately fitted update hazard with post-hoc clipping to $q\ge r$;
\item the proposed MHB model, fitted with coherence enforced throughout.
\end{enumerate}

Prices live on a half-cent mid-quote lattice. The primary log score therefore
uses the forecast probability of the observed half-cent cell. For MHB this cell
probability includes both integrated beta density and the no-update atom. For a
clipped Gaussian it includes the normal probability mapped into the cell and
the boundary mass created by clipping. This makes the score proper on a common
discrete observation space and avoids assigning an arbitrary density unit to
one model but not another.

We also report 50, 80, and 95 percent interval coverage, interval scores
\citep{winkler1972decision,gneiting2007strictly}, update-hazard Brier scores,
randomized probability integral transforms, and both equal- and volume-weighted
aggregates. Contract-cluster resampling preserves within-contract dependence.
All preprocessing thresholds and subgroup analyses are fixed before the full
test outcomes are evaluated.

\section{Public-API replication design}

The independent replication uses Kalshi's unauthenticated historical market
and candlestick endpoints. Hourly candlesticks contain closing YES bid and ask,
last price, volume, and open interest. The forecast price is the closing
mid-quote. Only consecutive observations exactly one hour apart are retained.
The scheduled expected-expiration field defines time to resolution; realized
settlement timing is never used as a forecast-origin feature.

The frozen panel spans August 2021 through April 2026 and includes prespecified
series from Economics, Climate and Weather, Politics, and Sports. Markets must
have at least 48 valid one-hour forecast origins and at least 100 contracts of total volume. The
main quote-quality rule requires a closing spread no wider than 20 cents;
thresholds of 5, 10, 50, and 100 cents are sensitivity checks. Raw JSON is
cached, every artifact is hashed, and failed requests are recorded in the run
manifest.

This panel is an independent replication, not a claim of row-level identity
with \citet{xi2026volatility}. Historical API aggregation, market metadata
updates, and quote cleaning may differ from the authors' construction. The
original headline ordering is treated as a reproduction target, but any
difference is reported rather than silently tuned away.

\section{Results}

\IfFileExists{generated/results.tex}{\subsection{Locked out-of-sample results}

The public-API acquisition selected 4,416 contracts and successfully retrieved 4,416. After forecast-validity, quote-quality, minimum-history, and expanding-window rules, the locked evaluation contains 1,007,560 hourly forecasts on 1,404 contracts across 50 test months. The unweighted update rate is 17.3\%; its forecast-origin volume-weighted counterpart is 44.5\%.

\begin{table}[ht]
\centering
\caption{Full-panel out-of-sample distribution forecasts. Lower is better for both scores.}
\label{tab:main-results}
\small
\begin{tabular}{lrrrr}
\toprule
Model & Tick log score & 95\% IS & Coverage & Width \\
\midrule
DR normal & 9.5984 & 0.2413 & 0.876 & 0.0971 \\
DR--AS normal & 3.1246 & 0.1896 & 0.980 & 0.1483 \\
Beta, no hurdle & 6.2704 & 0.2125 & 0.949 & 0.1286 \\
Hurdle beta, post-hoc & 6.0321 & 0.2148 & 0.949 & 0.1386 \\
MHB & 5.5271 & 0.2151 & 0.949 & 0.1390 \\
\bottomrule
\end{tabular}
\end{table}

On active updates, the independent replication preserves the source paper's structural ordering: DR--AS posts a volume-weighted 95\% interval score of 0.2744, versus 0.4412 for deadline resolution alone. The DR-minus-DR--AS improvement is 0.1667 with a 95\% contract-cluster interval [0.0158, 0.2620].

Volume is concentrated: the two largest contracts contribute 50.3\% of evaluation weight. Under the pre-specified contract-balanced robustness rule, MHB's tick log score is 6.5503, versus 4.0972 for DR--AS normal and 8.8790 for the post-hoc hurdle; the corresponding interval scores are 0.2250, 0.2167, and 0.2246. Under equal observation weights, by contrast, MHB has the best tick log score (1.9974 versus 3.2032 for DR--AS) but a worse interval score (0.0787 versus 0.0750).

\begin{table}[ht]
\centering
\caption{Volume-weighted category results (descriptive).}
\label{tab:category-results}
\small
\begin{tabular}{lrrrr}
\toprule
Category & DR--AS log & MHB log & DR--AS IS & MHB IS \\
\midrule
Economics & 3.8702 & 3.4998 & 0.1857 & 0.1933 \\
Politics & 2.9638 & 3.4579 & 0.1589 & 0.1934 \\
Sports & 3.1669 & 9.7104 & 0.2435 & 0.2593 \\
\bottomrule
\end{tabular}
\end{table}

The MHB density advantage over DR--AS appears in Economics but reverses in Politics and Sports; interval scores do not improve. These category results are descriptive, not multiplicity-adjusted. The two pre-specified daily weather series contribute no contracts after the 48-valid-origin rule, which is a limitation of this public-API replication rather than a basis for changing the locked filter.

Against the active-fit DR--AS normal benchmark, MHB's tick log score is higher by 2.4025 (95\% contract-cluster interval [0.8307, 6.4983]).
Against the active-fit DR--AS normal benchmark, MHB's 95\% interval score is higher by 0.0255 (95\% contract-cluster interval [0.0107, 0.0374]).

Against the beta law without a hurdle, MHB's tick log score is lower by 0.7433 (95\% contract-cluster interval [0.6201, 1.0664]).
Against the beta law without a hurdle, the beta law without a hurdle-minus-MHB 95\% interval score difference is -0.0026; the 95\% contract-cluster interval [-0.0145, 0.0045] includes zero.

Against the post-hoc coherent hurdle, MHB's tick log score is lower by 0.5049 (95\% contract-cluster interval [0.1329, 1.4202]).
Against the post-hoc coherent hurdle, the post-hoc coherent hurdle-minus-MHB 95\% interval score difference is -0.0004; the 95\% contract-cluster interval [-0.0011, 0.0000] includes zero.

\begin{table}[ht]
\centering
\caption{Volume-weighted scores by update regime.}
\label{tab:regime-results}
\small
\begin{tabular}{llrr}
\toprule
Regime & Model & Tick log score & 95\% IS \\
\midrule
All & DR--AS normal & 3.1246 & 0.1896 \\
 & Beta, no hurdle & 6.2704 & 0.2125 \\
 & Hurdle beta, post-hoc & 6.0321 & 0.2148 \\
 & MHB & 5.5271 & 0.2151 \\
\addlinespace
Active & DR--AS normal & 3.9144 & 0.2744 \\
 & Beta, no hurdle & 7.9759 & 0.3426 \\
 & Hurdle beta, post-hoc & 9.4564 & 0.3378 \\
 & MHB & 8.3244 & 0.3385 \\
\addlinespace
Inactive & DR--AS normal & 2.4925 & 0.1218 \\
 & Beta, no hurdle & 4.9054 & 0.1084 \\
 & Hurdle beta, post-hoc & 3.2915 & 0.1164 \\
 & MHB & 3.2884 & 0.1164 \\
\bottomrule
\end{tabular}
\end{table}

MHB central-interval coverage at the pre-specified levels is 50\%: 67.2\%, 80\%: 87.7\%, 95\%: 94.9\%. Its volume-weighted update-hazard Brier score is 0.2307, versus 0.2283 for the separately fitted hazard. The separate hazard violates $q\ge r$ at 5.49\% of test origins before post-hoc repair.

\begin{table}[ht]
\centering
\caption{Pre-specified spread-threshold sensitivity.}
\label{tab:spread-sensitivity}
\small
\begin{tabular}{lrrrr}
\toprule
Max spread & Post-hoc log & MHB log & Post-hoc IS & MHB IS \\
\midrule
0.05 & 6.1602 & 5.5615 & 0.2168 & 0.2172 \\
0.10 & 6.0445 & 5.5344 & 0.2151 & 0.2155 \\
0.20 & 6.0321 & 5.5271 & 0.2148 & 0.2151 \\
0.50 & 6.0268 & 5.5250 & 0.2149 & 0.2153 \\
1.00 & 6.0608 & 5.5246 & 0.2151 & 0.2155 \\
\bottomrule
\end{tabular}
\end{table}
}{
\begin{center}
\fbox{\parbox{0.88\linewidth}{\textbf{Locked-results placeholder.}
The full expanding-window run, cluster bootstrap, multi-level calibration, and
spread-threshold sensitivity analysis must finish before empirical claims are
inserted. The one-month smoke test is excluded from the manuscript evidence.}}
\end{center}
}

\section{Limitations and interpretation}

First, the displayed quote is not the frictionless efficient price. Exact
staleness can reflect no information, thin liquidity, exchange closure, or a
latent move smaller than one tick. The rounding-aware score addresses the price
lattice but does not identify these mechanisms separately.

Second, the active beta law is a compact forecasting kernel rather than an exact
transition density. It is useful because it is bounded, asymmetric, and moment
coherent. Applications that require exact absorption probabilities should use a
Wright-Fisher spectral or high-accuracy approximate transition and retain the
same coherence law.

Third, a known resolution deadline is not the same as a known information
event. Sports games, court decisions, and macro announcements concentrate
information at scheduled or state-dependent times. The current hazard absorbs
some of that predictability through time, activity, and lagged updates, but a
calendar of scheduled events is the natural next covariate set.

Finally, market prices equal risk-neutral rather than necessarily physical
probabilities. The martingale result is a pricing-measure statement. Risk
premia, fees, position limits, and manipulation can create predictable physical
drift even if discounted prices are martingales under a selected measure.

\section{Conclusion}

In a bounded prediction market, ``whether a price moves'' and ``how volatile it
is'' are not separable modeling choices. The martingale and binary-payoff
structure imposes the sharp inequality $r\le q$. Enforcing that law turns the
DR-AS variance insight into an unconditional predictive distribution with a
no-change atom, bounded support, and proper density forecasts. Whether the added
structure improves out-of-sample performance is an empirical question; this
repository is designed so that the answer can be negative without compromising
the reproducibility of the test.

\begin{thebibliography}{99}

\bibitem[Archak and Ipeirotis(2010)]{archak2010modeling}
N. Archak and P. G. Ipeirotis.
Modeling volatility in prediction markets.
In \emph{Proceedings of the ACM Conference on Electronic Commerce}, pages 197--206, 2010.

\bibitem[Catania et al.(2019)]{catania2019dynamic}
L. Catania, R. Di Mari, and P. Santucci de Magistris.
Dynamic discrete mixtures for high frequency prices.
SSRN working paper 3349118, 2019.

\bibitem[Dalen(2025)]{dalen2025toward}
S. Dalen.
Toward Black-Scholes for prediction markets: A unified kernel and market-maker's handbook.
arXiv:2510.15205, 2025.

\bibitem[Gneiting and Raftery(2007)]{gneiting2007strictly}
T. Gneiting and A. E. Raftery.
Strictly proper scoring rules, prediction, and estimation.
\emph{Journal of the American Statistical Association}, 102(477):359--378, 2007.

\bibitem[Hol\'y(2023)]{holy2023intraday}
V. Hol\'y.
An intraday GARCH model for discrete price changes and irregularly spaced observations.
arXiv:2211.12376, 2023.

\bibitem[Roa et al.(2022)]{roa2022smalltime}
T. Roa, M. I. Fariello, G. Mart\'inez, and J. Le\'on.
Small-time approximation of the transition density for diffusions with singularities: Application to the Wright-Fisher model.
arXiv:2212.11442, 2022.

\bibitem[Winkler(1972)]{winkler1972decision}
R. L. Winkler.
A decision-theoretic approach to interval estimation.
\emph{Journal of the American Statistical Association}, 67(337):187--191, 1972.

\bibitem[Xi et al.(2026)]{xi2026volatility}
W. Xi, C. C. Moallemi, M. Pai, and S. Wang.
Volatility in prediction markets: A structural approach.
arXiv:2607.08199v2, 2026.

\bibitem[Xu(2026)]{xu2026local}
L. Xu.
A local-volatility theory of prediction markets: Absorbed martingales on the simplex and an identification law for transient mispricing.
SSRN working paper 7012278, 2026.

\end{thebibliography}

\end{document}
